\chapter{Preliminarii}

	În acest capitol sunt stabilite noțiunile de bază privind analiza vibrațiilor pentru diagnosticarea motoarelor și principiile de funcționare ale rețelelor neuronale convoluționale unidimensionale (1D CNN). De asemenea, este prezentat și analizat stadiul actual al soluțiilor de mentenanță predictivă, cu accent pe integrarea paradigmei TinyML și a procesării prin Edge Computing, definind astfel contextul în care se raportează obiectivele acestei lucrări.
	
	Vibrațiile mecanice sunt un indicator important pentru identificarea stării de funcționare a unui sistem industrial. \cite{wowk1991machinery} Studiile arată că problemele mecanice incipiente pot genera anomalii vibraționale subtile  sau al unui tipar specific și pot fi detectate cu anumiți senzori de sunet sau mișcare. \cite{vitolo2022low, kolok2025low, bala2024artificial, garcia2025condition} Monitorizarea mișcării se realizează de regulă printr-un accelerometru, acesta oferă informații sensibile și detaliate despre  anomalii precum dezechilibru, dezaliniere sau uzură. \cite{gupta2025embedded} Astfel mediile industriale stricte au impus abordarea bazată pe inteligența artificială. \cite{vitolo2022low}
	
	Pentru dezvoltarea unui sistem eficient de mentenanță predictivă este necesară folosirea unor tehnologii precum: preluarea, stocarea și analiza semnalelor, rețele neuronale aprofundate și Edge Computing.
	
	Datorită tehnologiei TinyML (TinyML e definit ca inferență pe dispozitive cu putere scăzută și memorie RAM limitată) un sistem industrial poate face analiză la nivel local cu o rată de eșantionare ridicată a fluxului de date, producând alerte în timp real, permițând o mentenanță predictivă fără a transfera un volum mare de date în cloud. \cite{gupta2025embedded}
	
	Pentru transpunerea în realitate a arhitecturilor TinyML se folosesc platforme hardware în apropiere sau pe același cip cu senzorul. O soluție populară și extrem de eficientă este utilizarea microcontrolerelor precum ESP32, care dispun de procesoare performante, în cazul ESP32 de un procesor Dual Core, care permit și module de conectivitate wireless și bluetooth, dar care să mențină consumul energetic redus. \cite{vitolo2022low, gupta2025embedded} Acesta lucrează direct cu accelerometrul și alți senzori pentru achiziția semnalului. \cite{gupta2025embedded}
	
	
	Pentru implementarea algoritmilor pe microcontrolerul ESP32 s-a utilizat framework-ul TensorFlow Lite for Microcontrollers \cite{tensorflow:tflitemicro_github}, integrat în mediul de dezvoltare prin portarea dedicată TensorFlowLite\_ESP32 \cite{tanaka2026tflite_esp32}, o soluție open-source dezvoltată special pentru dispozitive cu memorie limitată. Această bibliotecă permite cuantizarea modelului și stocarea lui în memoria Flash, și executarea inferențelor în timp real pe semnalul brut cu un Tensor Arena alocat în SRAM. \cite{esp32_datasheet, tensorflow:tflitemicro_github} 
	
\section{Stadiul Actual al Subdomeniului}
	
	Soluțiile clasice de diagnosticarea defecțiunilor se bazează frecvent pe soluții cu infrastructură Cloud, care necesită o lățime de bandă mare, costuri ridicate, dependență permanentă de o conexiune și latență mare, unde se adaugă problema confidențialității și securizării datelor transmise. \cite{bala2024artificial, alfarag3embedded, gupta2025embedded} Edge Computing presupune mutarea procesării direct pe dispozitivul local, permițând reacții în timp real și funcționare independentă, de obicei cu costuri mai scăzute decât infrastructura cloud. În această paradigmă se regăsește tehnologia de TinyML care scalează și cuantizează algoritmii de machine learning, permițând astfel rularea algoritmilor de învățare automată direct pe dispozitive embedded cu o putere joasă de calcul și resurse limitate printre care și cele energetice. Pentru a permite astfel de implementări rețeaua neuronală trece prin tehnici specifice de optimizare precum: cuantizarea care reprezintă convertirea parametrilor rețelei din format cu virgulă mobila în numere întregi mai compacte și Model Pruning care elimină conexiunile redundante din rețea. \cite{gupta2025embedded, bala2024artificial}
	
	Deși rețelele neuronale convoluționale bidimensionale (2D CNN) au devenit un standard pentru recunoașterea imaginilor, ele sunt proiectate pentru date bidimensionale (grile de pixeli), iar semnalele furnizate de senzori sunt serii temporale unidimensionale. Astfel, aplicarea unui CNN 2D necesită transformarea prealabilă a semnalului într-o reprezentare bidimensională, ceea ce implică un pas suplimentar de procesare și cerințe mai mari de memorie și putere de calcul. \cite{kiranyaz20191, bitam2025artificial} De aceea, CNN-ul 1D este frecvent utilizat în procesarea semnalelor vibraționale în timp real. \cite{kiranyaz20191, bala2024artificial} Printre avantajele principale se numără posibilitatea de a opera direct pe semnalul brut \cite{kiranyaz20191, ince2016real}, fuziunea extragerii caracteristicilor diferite într-un singur proces de învățare \cite{kiranyaz20191, ince2016real} și complexitatea computațională redusă, deoarece efectuează convoluții strict unidimensionale, necesitând astfel un număr mic de neuroni ascunși și straturi. \cite{kiranyaz20191} Fiind o arhitectură compactă, permite rularea modelului direct pe platforme Edge. \cite{bala2024artificial, gupta2025embedded}
	
	Fezabilitatea rulării algoritmilor de învățare automată direct pe dispozitive edge pentru mentenanță predictivă a fost demonstrată în mai multe studii. \cite{gupta2025embedded, bala2024artificial} Totuși, evaluările comparative și neutre realizate direct pe hardware-ul de margine rămân relativ puține, studiile de sinteză recente în domeniul TinyML semnalând absența unor benchmark-uri robuste și standardizate. \cite{alharthi2026tinyml, somvanshi2025tiny} Un exemplu este studiul realizat de Alfarag (2026), care a antrenat mai mulți algoritmi clasici și rețele neuronale precum Random Forest, XGBoost, LSTM și CNN 1D, folosind seturi identice de date, pentru asigurarea evaluării neutre. \cite{alfarag3embedded} Aceste modele au fost comprimate și implementate fizic pe microcontrolere cu resurse limitate, precum Arduino Portenta și Raspberry Pi Pico. \cite{alfarag3embedded} Rezultatele testelor practice, realizate prin conectarea directă a senzorilor la microcontrolere, au arătat aplicabilitatea algoritmilor de machine learning la marginea rețelei, cu timpi de inferență reduși, sub 20 de milisecunde pentru CNN 1D. \cite{alfarag3embedded} Spre deosebire de acest studiu, care folosește platforme mai puternice și nu publică setul de date, lucrarea de față realizează evaluarea pe un ESP32 low-cost, cu un set de date propriu și timpi de inferență măsurați direct pe dispozitiv.
	
	Deși în mediile simulate au demonstrat succes în implementarea algoritmilor de învățare automată și funcționarea lor pentru a face mentenanță predictiva, implementările de tip end-to-end în lumea reală se lovesc frecvent de limitările fizice hardware, precum sensibilitatea și pragul de zgomot al senzorilor. De exemplu, senzorii MEMS comerciali utilizați la marginea rețelei precum MPU6050 prezintă un nivel de zgomot mai ridicat comparativ cu senzorii piezoelectrici folosiți în industrie, captând cu ușurință zgomotul electric sau interferențele electromagnetice generate de motor. Fără aplicarea unor tehnici de atenuare a zgomotului la nivel hardware și software, precum filtrul intern trece-jos al senzorului, rețeaua neuronală riscă să extragă caracteristici din zgomotul electric ambiental în loc să identifice anomalia mecanică reală. Astfel, filtrarea semnalului este o etapă critică a preprocesării, acesta asigurând echilibru între claritatea datelor și simplitatea computațională pentru a garanta că modelul evaluează doar componentele de joasă frecvență asociate cu problemele mecanice. \cite{kolok2025low}
	
